\begin{abstract}
针对塔式太阳能光热发电站定日镜场的优化设计问题，本文建立了一套基于光学几何原理的镜场性能仿真与优化模型。首先，构建了太阳位置与DNI计算模型以及包含余弦效率、阴影遮挡效率、大气透射率和截断效率的光学效率综合模型。通过将镜面离散化为网格点并采用光线追踪思想，精确计算了每面定日镜在各时点的阴影遮挡效率。基于该仿真模型，进一步构建了以“单位镜面面积年平均输出热功率”最大化为目标的优化模型，采用径向交错布局策略与遗传算法进行协同优化。问题1仿真结果：在1745面定日镜、总镜面面积62820.0 m²下，年平均光学效率为0.6699，年平均输出热功率为30.0887 MW，单位面积功率为0.479 kW/m²，未达到60 MW额定功率。各月平均功率呈现季节性变化，6月峰值36.1031 MW，12月谷值21.5355 MW。问题2与问题3的优化模型以60 MW额定功率为约束，通过优化吸收塔位置、定日镜尺寸、安装高度及布局参数，旨在提升单位镜面面积功率至0.479 kW/m²以上。灵敏度分析表明，模型对镜面反射率和镜面宽度高度敏感（变化±10%时输出功率变化约±10%），而对安装高度不敏感（变化±20%时输出功率变化小于0.1%）。本文所提出的模型与方法，为塔式光热电站的定日镜场高效设计提供了理论依据与数值参考。

\keywords{塔式太阳能光热发电；定日镜场优化；光学效率；遗传算法；单位镜面面积年平均输出热功率；灵敏度分析；径向交错布局}
\end{abstract}